% Generated draft: requires mathematical and source review.
\documentclass[11pt]{article}
\usepackage[margin=25mm]{geometry}
\usepackage{fontspec}
\setmainfont{Latin Modern Roman}
\usepackage{amsmath,amssymb,mathtools,unicode-math}
\setmathfont{Latin Modern Math}
\usepackage{enumitem,needspace,xurl,hyperref}
\setlength{\parindent}{0pt}
\setlength{\parskip}{5pt}
\setlength{\emergencystretch}{3em}
\newcommand{\N}{\mathbb{N}}
\newcommand{\Z}{\mathbb{Z}}
\newcommand{\Q}{\mathbb{Q}}
\newcommand{\R}{\mathbb{R}}
\newcommand{\C}{\mathbb{C}}
\newcommand{\F}{\mathbb{F}}
\newcommand{\A}{\mathbb{A}}
\newcommand{\PP}{\mathbb{P}}
\newcommand{\E}{\mathbb{E}}
\newcommand{\eps}{\varepsilon}
\newcommand{\dd}{\,\mathrm{d}}
\begin{document}
\section*{UnsolvedMath ID: 30002512}
\textbf{Source problem code:} catalogue entry
\subsection{Definitions and mathematical statement}
For reference, two smooth Riemannian metrics \(g\) and \(\widehat g\) on the same smooth surface \(S\) are classically conformally equivalent if there exists \(u\in C^\infty(S,\R)\) such that
\[
\widehat g=e^{2u}g.
\]
This standard definition provides context; the supplied record does not specify whether its intended convergence concerns metrics, conformal structures, conformal maps, or uniformizations.

\textbf{Source question (2014).} Prove or disprove that discrete conformality converges to classical conformality as the triangulations of a surface become arbitrarily fine.

\textbf{Unresolved formulation.} The supplied record does not define discrete conformality, the discrete data associated with a triangulation, or their comparison with classical conformal data. It does not specify the class of surfaces, including compactness, orientability, genus, or the presence of boundary; the regularity and geometric data on the surface; the meaning of refinement or arbitrarily small mesh; admissible mesh-quality conditions; a topology of convergence; or any marking, normalization, or boundary conditions. It also does not specify whether convergence is required for every admissible sequence of triangulations or for a suitably chosen sequence. Consequently, the source question cannot be converted into a uniquely determined quantified convergence assertion from the supplied evidence alone. These choices remain unresolved rather than being imposed as additional hypotheses.

The supplied dated assessment labels the problem partially solved and reports convergence results for closed surfaces under mesh-quality hypotheses, with separate positive results in positive genus and genus zero, while describing the unrestricted conjectural formulation as not fully settled. This is supplied curation, not an independently verified mathematical or literature conclusion.
\subsection{Short English statement}
The question asks whether discrete conformality approaches classical conformality as surface triangulations become arbitrarily fine. A precise convergence claim requires definitions of the discrete model, admissible surfaces and triangulations, comparison objects, and convergence topology. These are absent from the supplied record. Its partially solved label and reported restricted results are supplied dated assessments requiring independent review.
\subsection{Sources}
\begin{itemize}
\item Local registry: \texttt{data/unsolvedmath/problems.json}, permanent UnsolvedMath ID \(30002512\), source problem code \texttt{catalogue entry}, title ``Convergence of Discrete to Classical Conformality.'' The supplied record assigns proposed year 2014 and names no individual proposer.
\item Supplied report citation: \emph{New Trends in Teichm\"uller Theory and Mapping Class Groups} (2014), Oberwolfach Reports, pp.\ 395--452, DOI 10.4171/owr/2014/07. Listed literature source: ``Oberwolfach Report 07/2014,'' \url{https://doi.org/10.4171/owr/2014/07}. The supplied background describes Oberwolfach Reports as documenting mathematical workshops and their open-problem discussions; no problem-specific page or passage is supplied.
\item Supplied literature source: \emph{Discrete uniformizations for closed surfaces}, \url{https://arxiv.org/abs/2008.06744}.
\item Supplied literature source: \emph{Discrete conformality in genus zero}, \url{https://arxiv.org/abs/2110.08208}.
\item Literature assessment date: \texttt{2026-08-21}. The supplied assessment reports positive convergence results for closed surfaces under mesh-quality hypotheses in positive genus and genus zero, and describes the unrestricted formulation as not fully settled. This assessment is supplied dated curation derived from the OWR material, not independently verified here; its time-sensitive status and cited claims require independent verification.
\item The supplied registry labels the statement \texttt{corrected\_verified} and says its mathematical content was unchanged while notation, typography, and wording were normalized. These are supplied provenance claims and do not establish independent verification of this draft.
\end{itemize}
\end{document}
